\section{The exact price by rule class}
\label{sec:exact}

The six rule groups fall into four classes by how a boost on one task moves the target's standing (\Cref{tab:rules}). Within a class the cheapest plan has a closed form or a small exact program; this lets us report optima, not a heuristic's, on boards with thousands of models. For rules whose rival scores do not depend on the target (mean family, quantiles), let $V=\max_{i\ne t}A_i(S)+\varepsilon$ be the level the target's aggregate must reach, and $d_j=V-S_{tj}$ the shortfall on task $j$ in the case of a quantile rule. Costs are linear in points unless stated.

\begin{table}[t]
\caption{Rule classes, what a boost buys, and the price each class implies. $a$ is the number of tasks on which the target already scores at least the rival level $V$, $r$ the rank of the order statistic ($r=m$ for the minimum, $r=\lceil m/2\rceil$ for the median), $G$ the aggregate gap, $u_j=\min(B,c_j-S_{tj})$ the headroom on task $j$.}
\label{tab:rules}
\centering\small
\begin{tabular}{@{}p{2.3cm}p{2.7cm}p{3.1cm}p{3.0cm}@{}}
\toprule
Rule & Influence of a boost & Minimal breadth & Depth at that breadth\\
\midrule
Mean family & linear, capped by headroom & fewest tasks whose headroom $\sum w_ju_j$ covers $G$ & fixed: $mG$ (raw mean), however spread\\
Quantile, median, min & none beyond level $V$ & $\max(0,r-a)$ & the $r-a$ smallest shortfalls $d_j$\\
Mean rank (= mean win rate) & steps, bounded per task & at least $\lceil g/N\rceil$ for deficit $g$ & optimum of a covering program (levels per task)\\
Plurality, Copeland & steps, coupled across tasks & optimum of an integer program & its optimum\\
\bottomrule
\end{tabular}
\end{table}

\begin{proposition}[Quantile rules]\label{prop:quantile}
Let $A$ be the $r$-th largest task score. The target needs at least $r$ tasks at level $V$ or above, so its minimal breadth is $\max(0,r-a)$. The cheapest plan with that breadth raises the $r-a$ tasks of smallest shortfall exactly to $V$; it is feasible under budget $B$ if and only if the $(r-a)$-th smallest shortfall is at most $B$ (and at most the task's headroom). Any boost beyond $V$ on a task, and any additional task, adds depth and removes nothing from the price.
\end{proposition}

\begin{proposition}[Mean family]\label{prop:mean}
Let $A_i=\sum_j w_jS_{ij}$ with weights $w_j>0$ (raw mean: $w_j=1/m$; $z$-scored mean: $w_j=1/(m\sigma_j)$) and gap $G=V-A_t>0$. First place is feasible under $B$ if and only if $\sum_jw_ju_j\ge G$. Under a linear cost $\sum_j\kappa_j\delta_j$ the optimum fills tasks in decreasing order of $w_j/\kappa_j$; for the raw mean in points the depth is $mG$ however it is spread over tasks. The minimal breadth is the least $k$ for which the $k$ largest $w_ju_j$ sum to at least $G$. Only under a scaled or convex cost does the allocation change the depth, and then it is water-filling.
\end{proposition}

\begin{proposition}[Additive rank rules]\label{prop:rank}
Under mean rank (Borda; with ties counted as half it equals the mean win rate) the margin of the target over a fixed rival $i$ is a sum over tasks of non-decreasing step functions of $\delta_j$, with at most $N-1$ steps per task (one for each model passed). The price is therefore the optimum of a multiple-choice covering problem with one constraint for each rival, and a boost beyond the best rival's score on a task buys nothing. In pairwise-win units each task can raise the margin over $i$ by at most $N$, so the breadth is at least $\lceil g_i/N\rceil$ for a deficit $g_i$.
\end{proposition}

\begin{proposition}[Coupled rules]\label{prop:coupled}
Under plurality the target wins a task only by passing the best rival there, and each such win also removes one win from that rival; under Copeland the target beats a rival only if it passes them on more tasks than they pass it. In both cases the effect of boosting a task depends on which other tasks are boosted, and the price is an integer program with one constraint per rival.
\end{proposition}

Proofs are in \Cref{app:proofs}.

\paragraph{Solvers and checks.}
We compute the optimum for the quantile and mean classes in closed form, and for the rank and coupled classes with mixed-integer programs (HiGHS) over the per-task levels, with rival constraints added lazily where a rule has one per rival; each plan returned is re-evaluated and is first under the rule on every price we report. Every solver agreed with exhaustive enumeration of integer plans on random instances of at most 7 models and 5 tasks (4 for the rank and coupled rules) (217 to 2072 instances per rule, all agreeing). Plurality is solved to optimality on every board. The mean rank and Copeland programs are solved to a reported optimality gap, which is zero at benchmark level and for the top-30 HELM targets at ceilings and up to 12\% (Borda) and 7\% (Copeland) on the top-30 sub-task targets; where it is not zero the price is an upper bound (\Cref{tab:price}).

\paragraph{No rule is cheapest everywhere.}
The classes give no universal ordering: with two models and three tasks, the target $(59,59,0)$ against a rival $(60,60,60)$ needs about 2 points under the median but about 62 under the mean (\Cref{app:proofs}). Whether and where this matters on real boards is the question of \Cref{sec:results}.
